\documentclass[../../../../main.tex]{subfiles}

\begin{document}

\section*{1989年 東京大学 理科}

\subsection*{第1問}

$k > 0$. $xy$ 平面上の 2 曲線 $y = k(x - x^3)$, $x = k(y - y^3)$ が第一象限に相異なる $(\alpha, \beta)$ を持つような $k$ の値域.

\bigskip

[解]
\[
\begin{cases}
l_1 : y = k(x - x^3) & \dots \text{①} \\
l_2 : x = k(y - y^3) & \dots \text{②}
\end{cases}
\]
とする. $(\alpha, \beta)$ は ① $\land$ ② をみたす. $(x, y)$ のまま変形して, ① $+$ ②, ① $-$ ② から,
\[
\begin{cases}
x + y = k \{(x + y) - (x^3 + y^3)\} \\
y - x = k \{(x - y) - (x^3 - y^3)\}
\end{cases}
\]
\[
\iff
\begin{cases}
x + y = k (x + y) \{1 - (x^2 - xy + y^2)\} \\
y - x = k (x - y) \{1 - (x^2 + xy + y^2)\} & \dots \text{③}
\end{cases}
\]
さて, 今交点 $(\alpha, \beta)$ は ③ をみたす. $\alpha \neq \beta$ 及び $0 < \alpha, \beta$ から
\[
\alpha + \beta > 0, \quad \alpha - \beta \neq 0
\]
だから, ③を $(x, y)$ のまま変形すれば,
\[
\begin{cases}
1 = k \{1 - (x^2 - xy + y^2)\} \\
-1 = k \{1 - (x^2 + xy + y^2)\}
\end{cases}
\]
\[
\iff
\begin{cases}
0 = 2k (1 - x^2 - y^2) \\
2 = 2k xy
\end{cases}
\]
\[
\iff
\begin{cases}
1 = x^2 + y^2 \\
\frac{1}{k} = xy \quad (\because k > 0) & \dots \text{④}
\end{cases}
\]
ゆえに, ④から「$0 < xy \dots \text{⑤}$」かつ「$x \neq y \dots \text{⑥}$」なる $(x, y)$ の存在条件をもとめれば良い. これは ④, ⑤, ⑥ のグラフを考えて, 及び $0 < k$ ゆえ $0 < \frac{1}{k}$ だから,
\[
0 < \frac{1}{k} < \frac{1}{2} \quad \therefore 2 < k
\]

\bigskip

[解注] 最後, グラフがあらだというのであれば,

\medskip

ア ④, ⑤, ⑥より, $x = \cos\theta, y = \sin\theta$ とおける ($0 < \theta < \pi/2, \theta \neq \pi/4$). ゆえに
\[
xy = \cos\theta \sin\theta = \frac{1}{2} \sin 2\theta \quad (0 < 2\theta < \pi, 2\theta \neq \pi/2)
\]
・ゆえに, この値域から,
\[
0 < xy < \frac{1}{2}
\]
を得る.

\medskip

イ ⑤から, ④において
\[
y = \frac{1}{k x}
\]
である. ④, ⑤, ⑥に代入して, $y$ を消去する.
\[
\begin{cases}
1 = x^2 + \frac{1}{k^2 x^2} \\
0 < x \\
x \neq \frac{1}{k x}
\end{cases}
\]
\[
\iff t^2 - t + \frac{1}{k^2} = 0 \quad (t = x^2) \dots \star
\]
\[
\begin{cases}
t > 0 \\
t \neq 1/k
\end{cases}
\]
ゆえにこれが解を持つ条件を調べれば良い. $\star$ の左辺 $f(t)$ とすれば
\[
f(1/k) = 1/k^2 > 0
\]
だから, $0 < t$ に 2 解持つことが必要十分で,
\[
D > 0 \land 0 < \frac{1}{2}
\]
\[
\iff 1 - 4/k^2 > 0
\]
\[
\iff 2 < k \quad (\because k > 0)
\]
をうる.

\subsection*{第2問}

[解] 題意の放物線は $y = \frac{1}{4} x^2$ である.

この上の点 $(a, \frac{1}{4} a^2)$ を中心とし, $y = -1$ に接する, つまり半径 $r$ として $r = 1 + \frac{1}{4} a^2$ の円 $C$ は
\[
C : (x - a)^2 + \left(y - \frac{1}{4} a^2\right)^2 = \left(1 + \frac{1}{4} a^2\right)^2
\]
である. $GF$ の中点 $I$ とすると, $PI \perp GF$ になっていることに注意して,
\[
T(a) = \frac{1}{2} \cdot IP \cdot GF = \frac{1}{2} \cdot IP \cdot (2 FI) \quad \dots \text{①}
\]

\begin{center}
\begin{tikzpicture}[scale=1.2, >=stealth]
  % Axes
  \draw[->] (-1.5,0) -- (4,0) node[right] {$y$};
  \draw[->] (0,-1.5) -- (0,4) node[above] {$z$};
  \node[below left] at (0,0) {$O$};

  % Line y = -1
  \draw[dashed] (-1.5,-1) -- (4,-1);
  \node[left] at (0,-1) {$-1$};

  % Parabola y = 1/4 x^2
  \draw[domain=-1.2:3.2, smooth, variable=\x] plot ({\x}, {0.25*\x*\x});

  % Points
  \coordinate (P) at (2.4, 1.44);
  \coordinate (F) at (0, 1);
  \coordinate (H) at (2.4, -1);

  \fill (P) circle (1.5pt) node[above right] {$P$};
  \fill (F) circle (1.5pt) node[left] {$F$};
  \fill (H) circle (1.5pt) node[below] {$H$};

  % Circle C
  \draw (P) circle (2.44);

  % Segments
  \draw (P) -- (H);
  \draw (P) -- (F);

  % Angle theta
  \draw (2.4,0.94) arc (-90:-168.6:0.5);
  \node at (2.1, 0.7) {$\theta$};

  % Point I
  \coordinate (I) at ($(F)!0.5!(P)$);
  \fill (I) circle (1.2pt) node[below right] {$I$};

  % Additional points G, C, D
  \node[above left] at (0.2, 3.8) {$G$};
  \node[above] at (1.5, 3.8) {$C$};
  \node[above right] at (3.2, 3.8) {$D$};

  % Chord GF perpendicular to PI
  \draw (F) -- (-0.2, 3.5);
\end{tikzpicture}
\end{center}

又, $PH$ と $PF$ のなす角 $\theta$ とすると, $\vec{PH} = \begin{pmatrix} 0 \\ -r \end{pmatrix}$, $\vec{PF} = \begin{pmatrix} -a \\ 1 - \frac{1}{4} a^2 \end{pmatrix}$ より,
\[
\cos\theta = \frac{\vec{PH} \cdot \vec{PF}}{|\vec{PH}| |\vec{PF}|} = \frac{-r \left(1 - \frac{1}{4} a^2\right)}{r^2} = \frac{-1 + \frac{1}{4} a^2}{r} = \frac{-4 + a^2}{4 + a^2} \quad \dots \text{②}
\]
この時, $\angle IPF = \frac{\pi}{2} - \theta$ にも注意して ①より
\[
T(a) = r \sin\left(\frac{\pi}{2} - \theta\right) r \cos\left(\frac{\pi}{2} - \theta\right) = r^2 \sin\theta \cos\theta
\]
\[
S(a) = \frac{1}{2} r^2 \theta
\]
だから,
\[
\frac{T(a)}{S(a)} = \frac{r^2 \sin\theta \cos\theta}{\frac{1}{2} r^2 \theta} = \frac{2 \sin\theta \cos\theta}{\theta} \quad \dots \text{③}
\]
ここで, ②より, $a \to \infty$ の時, $\cos\theta = \frac{-4/a^2 + 1}{4/a^2 + 1} \to 1$ より $\theta \to 0$ であるので, ③から,
\[
\frac{T(a)}{S(a)} = \frac{\sin 2\theta}{2\theta} \cdot 2 \to 2 \quad (a \to \infty)
\]
である.

\subsection*{第3問}

\[
H = \{ z = x + y i \mid x, y \in \mathbb{R}, y > 0 \}. \text{ 以下, } z \text{ を } H \text{ に属する複素数とする. } q \in \mathbb{R}_{>0} \text{ とし, } f(z) = \frac{z+1-q}{z+1} \text{ とおく.}
\]
\begin{enumerate}
  \item[(1)] $f(z) \in H$ を示せ.
  \item[(2)] $f_1(z) = f(z)$ とおき, $n = 2, 3, \dots$ に対し $f_n(z) = f(f_{n-1}(z))$ とおく. 全ての $H$ の元 $z$ に対し, $f_{10}(z) = f_1(z)$ が成立するような $q$ は?
\end{enumerate}

\bigskip

$\triangleright$ "$x + y i$" 形だから基本コレ

(1) では, $f(z)$ の虚部が正を示せば OK
\[
\operatorname{Im} f(z) = \frac{1}{2i} \{ f(z) - \overline{f(z)} \} > 0
\]

(2) は関数の問題. $f_{10} = f_1$ となる条件をもとめれば良いが, さすがに計はキツイから, $f_9 = z$ となる条件をもとめるか, 逆関数に注目.

\end{document}
